% Part: first-order-logic
% Chapter: syntax-and-semantics
% Section: satisfaction

\documentclass[../../../include/open-logic-section]{subfiles}

\begin{document}

\olfileid{fol}{syn}{sat}

\olsection{\article{structure} \printtoken{S}{structure}-এ
  \article{formula} \printtoken{S}{formula}-এর পরিতৃপ্তি}

\begin{explain}
এক দিকে পদ ও !!{formula}s-এর মতো রাশি এবং অন্য দিকে !!{structure}s-এর
মধ্যে যোগসূত্র স্থাপনকারী মৌলিক ধারণা হলো পদের \emph{!!{value}} এবং
!!a{formula}-এর \emph{পরিতৃপ্তি}। অনানুষ্ঠানিকভাবে, পদের !!{value}
হলো কোনো !!a{structure}-এর একটি !!{element}—পদটি শুধু একটি ধ্রুবক হলে
তার !!{value} হলো !!{structure}-টি ধ্রুবকটিকে যে বস্তু আরোপ করে, আর
পদটি !!{function}s দিয়ে গঠিত হলে তার !!{value}, পদের ধ্রুবকগুলির
!!{value}s এবং ফলনসংকেতগুলিকে আরোপিত ফলনগুলি থেকে গণনা করা হয়। কোনো
!!^a{formula} কোনো !!a{structure}-এ \emph{পরিতৃপ্ত} হয়, যদি
বিধেয়গুলিকে দেওয়া ব্যাখ্যায় !!{structure}-টির সংজ্ঞাক্ষেত্রে
!!{formula}-টি সত্য হয়। পরিতৃপ্তির এই ধারণাটি আরোহীভাবে নির্দিষ্ট:
!!{structure}-এর বিবরণ সরাসরি বলে দেয় কখন পরমাণু !!{formula}s
পরিতৃপ্ত, আর কোনো জটিল !!{formula} কখন পরিতৃপ্ত তা তার প্রধান সংযোজক
বা পরিমাণসূচক এবং তার অব্যবহিত !!{subformula}s পরিতৃপ্ত কি না, তার
ভিত্তিতে সংজ্ঞায়িত করি।

এখানে পরিমাণসূচকের ক্ষেত্রটি একটু কঠিন, কারণ পরিমাণিত কোনো
!!{formula}-এর অব্যবহিত !!{subformula}-তে মুক্ত !!{variable} থাকে,
কিন্তু !!{structure}s, !!{variable}s-এর !!{value}s নির্দিষ্ট করে না।
এই অসুবিধা সামলাতে আমরা \emph{চলরাশি-আরোপ} প্রবর্তন করি এবং পরিতৃপ্তিকে
শুধু কোনো !!a{structure}-এর সাপেক্ষে নয়, বরং কোনো !!a{structure} এবং
!!a{variable} আরোপ উভয়ের সাপেক্ষে সংজ্ঞায়িত করি।
\end{explain}

\begin{defn}[চলরাশি-আরোপ]
!!a{structure}~$\Struct{M}$-এর জন্য \emph{চলরাশি-আরোপ}~$s$ হলো এমন
একটি ফলন, যা প্রত্যেক !!{variable}-কে~$\Domain M$-এর কোনো
!!a{element}-এ চিত্রিত করে; অর্থাৎ $s\colon \Var \to \Domain M$।
\end{defn}

\begin{explain}
!!^a{structure} প্রত্যেক !!{constant}-কে !!a{value} আরোপ করে, আর
চলরাশি-আরোপ প্রত্যেক চলরাশিকে। কিন্তু এগুলি থেকে গঠিত পদগুলিকেও আমরা
!!{domain}-এর !!{element}s-এর নাম হিসেবে ব্যবহার করতে চাই। সে জন্য
পদের !!{value} আরোহীভাবে সংজ্ঞায়িত করি। !!{constant}s ও চলরাশির মান
!!{structure} বা চলরাশি-আরোপের নির্দিষ্ট করা মানই; আরও জটিল পদের মান
পুনরাবৃত্তিমূলকভাবে গণনা করা হয়, !!{structure} যে ফলনগুলি
!!{function}s-কে আরোপ করে সেগুলি ব্যবহার করে।
\end{explain}

\begin{defn}[পদের !!^{value}]
$t$ যদি $\Lang L$ ভাষার একটি পদ, $\Struct M$ যদি $\Lang L$-এর একটি
!!{structure} এবং $s$ যদি $\Struct M$-এর জন্য !!a{variable} আরোপ হয়,
তাহলে \emph{!!{value}}~$\Value{t}{M}[s]$ নিচের নিয়মে সংজ্ঞায়িত:
\begin{enumerate}
\item \indcase{t}{c}{$\Value{\indfrm}{M}[s] = \Assign{\indcomplex}{M}$।}
\item \indcase{t}{x}{$\Value{\indfrm}{M}[s] = s(\indcomplex)$।}
\item \indcase{t}{\Atom{f}{t_1, \ldots, t_n}}{
\[
\Value{\indfrm}{M}[s] = \Assign{f}{M}(\Value{t_1}{M}[s], \ldots,
\Value{t_n}{M}[s]).
\]}
\end{enumerate}
\end{defn}

\begin{defn}[$x$-বিকল্প]
$s$ যদি !!a{structure}~$\Struct M$-এর জন্য !!a{variable} আরোপ হয়,
তাহলে $\Struct M$-এর জন্য এমন যেকোনো !!{variable} আরোপ~$s'$-কে
$s$-এর \emph{$x$-বিকল্প} বলে, যা $x$-কে করা আরোপে সর্বাধিক $s$-এর
থেকে আলাদা। $s'$ যদি $s$-এর একটি $x$-বিকল্প হয়, লিখি
$\varAssign{s'}{s}{x}$।
\end{defn}

\begin{explain}
লক্ষ করুন, কোনো আরোপ~$s$-এর $x$-বিকল্পকে $x$-এর জন্য \emph{অবশ্যই}
আলাদা কিছু আরোপ করতে হয় না। বস্তুত প্রত্যেক আরোপই নিজের একটি
$x$-বিকল্প।
\end{explain}

\begin{defn}
  $s$ যদি !!a{structure}~$\Struct M$-এর জন্য !!a{variable} আরোপ হয়
  এবং $m \in \Domain{M}$, তাহলে আরোপ~$\Subst{s}{m}{x}$ হলো নিচের
  নিয়মে সংজ্ঞায়িত চলরাশি-আরোপ:
  \[\Subst{s}{m}{x}(y) = \begin{cases}
    m & \text{যদি } y \ident x\\
    s(y) & \text{অন্যথায়}.
  \end{cases}\]
\end{defn}

অন্যভাবে বললে, $\Subst{s}{m}{x}$ হলো $s$-এর সেই বিশেষ $x$-বিকল্প,
যা সংজ্ঞাক্ষেত্রের !!{element}~$m$-কে $x$-এর মান আরোপ করে এবং $x$
ছাড়া অন্য !!{variable}s-কে $s$-এর আরোপিত বস্তুগুলিই আরোপ করে।

\begin{defn}[পরিতৃপ্তি]
\ollabel{defn:satisfaction}
!!a{structure}~$\Struct M$-এ !!a{variable} আরোপ~$s$-এর সাপেক্ষে
!!a{formula}~$!A$-এর পরিতৃপ্তি, সংকেতে $\Sat{M}{!A}[s]$, নিচের নিয়মে
পুনরাবৃত্তিমূলকভাবে সংজ্ঞায়িত। (“$\Sat{M}{!A}[s]$ নয়” বোঝাতে আমরা
$\Sat/{M}{!A}[s]$ লিখি।)
\begin{enumerate}
\tagitem{prvFalse}{%
  \indcase{!A}{\lfalse}{$\Sat/{M}{\indfrm}[s]$।}}{}

\tagitem{prvTrue}{%
  \indcase{!A}{\ltrue}{$\Sat{M}{\indfrm}[s]$।}}{}

\item \indcase{!A}{\Atom{R}{t_1, \dots, t_n}}{$\Sat{M}{\indfrm}[s]$
  তখনই, যখন $\langle \Value{t_1}{M}[s], \dots, \Value{t_n}{M}[s]
  \rangle \in \Assign{R}{M}$।}

\item \indcase{!A}{\eq[t_1][t_2]}{$\Sat{M}{\indfrm}[s]$ তখনই, যখন
  $\Value{t_1}{M}[s] = \Value{t_2}{M}[s]$।}

\tagitem{prvNot}{%
  \indcase{!A}{\lnot !B}{$\Sat{M}{\indfrm}[s]$ তখনই, যখন
    $\Sat/{M}{!B}[s]$।}}{}

\tagitem{prvAnd}{%
  \indcase{!A}{(!B \land !C)}{$\Sat{M}{\indfrm}[s]$ তখনই, যখন
    $\Sat{M}{!B}[s]$ এবং $\Sat{M}{!C}[s]$।}}{}

\tagitem{prvOr}{%
  \indcase{!A}{(!B \lor !C)}{$\Sat{M}{\indfrm}[s]$ তখনই, যখন
    $\Sat{M}{!B}[s]$ অথবা $\Sat{M}{!C}[s]$ (অথবা উভয়ই)।}}{}

\tagitem{prvIf}{%
  \indcase{!A}{(!B \lif !C)}{$\Sat{M}{\indfrm}[s]$ তখনই, যখন
    $\Sat/{M}{!B}[s]$ অথবা $\Sat{M}{!C}[s]$ (অথবা উভয়ই)।}}{}

\tagitem{prvIff}{%
  \indcase{!A}{(!B \liff !C)}{$\Sat{M}{\indfrm}[s]$ তখনই, যখন হয়
    $\Sat{M}{!B}[s]$ ও $\Sat{M}{!C}[s]$ উভয়ই, নয়তো
    $\Sat{M}{!B}[s]$ ও $\Sat{M}{!C}[s]$ কোনোটিই নয়।}}{}

\tagitem{prvAll}{%
  \indcase{!A}{\lforall[x][!B]}{$\Sat{M}{\indfrm}[s]$ তখনই, যখন
    প্রত্যেক !!{element}~$m \in \Domain M$-এর জন্য
    $\Sat{M}{!B}[\Subst{s}{m}{x}]$।}}{}

\tagitem{prvEx}{%
  \indcase{!A}{\lexists[x][!B]}{$\Sat{M}{\indfrm}[s]$ তখনই, যখন অন্তত
  একটি !!{element}~$m \in \Domain M$-এর জন্য
  $\Sat{M}{!B}[\Subst{s}{m}{x}]$।}}{}
\end{enumerate}
\end{defn}

\begin{explain}
শেষ
\iftag{notprvEx,notprvAll}{অনুচ্ছেদটিতে}{দুটি অনুচ্ছেদে} চলরাশি-আরোপ
গুরুত্বপূর্ণ। \iftag{prvAll}{“সব $m \in \Domain{M}$-এর জন্য
$\Sat{M}{!B(m)}$”—এভাবে $\lforall[x][!B(x)]$-এর পরিতৃপ্তি সংজ্ঞায়িত
করতে পারি না।}{}\iftag{prvEx}{“অন্তত একটি $m \in \Domain{M}$-এর জন্য
$\Sat{M}{!B(m)}$”—এভাবে $\lexists[x][!B(x)]$-এর পরিতৃপ্তি সংজ্ঞায়িত
করতে পারি না।}{} কারণ $m \in \Domain M$ হলে $m$ ভাষার সংকেত নয়,
তাই $!B(m)$ !!a{formula} নয় (অর্থাৎ $\Subst{!B}{m}{x}$ অসংজ্ঞায়িত)।
আবার $\Struct{M}$-এর প্রত্যেক !!{element}-এর নাম দেয় এমন যথেষ্ট
!!{constant}s বা পদ আমাদের আছে ধরে নিতেও পারি না, কারণ
!!{structure}s-এর সংজ্ঞায় তার কোনো শর্ত নেই। মানক ভাষায়
!!{constant}s-এর সেট !!{denumerable}; তাই $\Domain{M}$
!!{enumerable} না হলে প্রত্যেক বস্তুর নাম দেওয়ার মতো যথেষ্ট
!!{constant}s-ই নেই।

চলরাশিকে সংজ্ঞাক্ষেত্রের !!{element}s-এর সঙ্গে সরাসরি যুক্ত করতে পারে
এমন !!{variable} আরোপ প্রবর্তন করে সমস্যাটির সমাধান করি। তখন, যেমন,
\iftag{prvEx}{$\lexists[x][!B(x)]$}{$\lforall[x][!B(x)]$},~$\Struct M$-এ
পরিতৃপ্ত তখনই যখন \iftag{prvEx}{অন্তত একটি $m \in
\Domain{M}$-এর জন্য}{সব $m \in \Domain{M}$-এর জন্য
$\Sat{M}{!B(m)}$}—এ কথা না বলে বলি, এটি~$\Struct M$-এ $s$-এর
\emph{সাপেক্ষে} পরিতৃপ্ত তখনই, যখন \iftag{prvEx}{অন্তত একটি}{প্রত্যেক}
$m \in \Domain M$-এর জন্য~$\Subst{s}{m}{x}$-এর সাপেক্ষে $!B(x)$
পরিতৃপ্ত।
\end{explain}

\begin{ex}
ধরি $\Lang{L} = \{a, b, f, R\}$, যেখানে $a$ ও $b$ হলো
!!{constant}s, $f$ একটি দ্বি-স্থানীয় !!{function} এবং $R$ একটি
দ্বি-স্থানীয় !!{predicate}। নিচের নিয়মে সংজ্ঞায়িত
!!{structure}~$\Struct{M}$ বিবেচনা করুন:
\begin{enumerate}
\item $\Domain M = \{1, 2, 3, 4\}$
\item $\Assign{a}{M} = 1$
\item $\Assign{b}{M} = 2$
\item $x+y \le 3$ হলে $\Assign{f}{M}(x, y) = x+y$, অন্যথায় $= 3$।
\item $\Assign{R}{M} = \{\tuple{1, 1}, \tuple{1, 2}, \tuple{2, 3}, \tuple{2, 4}\}$
\end{enumerate}
প্রত্যেক !!{variable}-কে $1 \in \Domain{M}$ আরোপকারী ফলন $s(x) = 1$
হলো~$\Struct{M}$-এর একটি চলরাশি-আরোপ।

তাহলে
\begin{align*}
\Value{f(a,b)}{M}[s] & = \Assign{f}{M}(\Value{a}{M}[s], \Value{b}{M}[s]).
\intertext{$a$ ও $b$ হলো !!{constant}s বলে $\Value{a}{M}[s]
  = \Assign{a}{M} = 1$ এবং $\Value{b}{M}[s] = \Assign{b}{M} = 2$। তাই}
\Value{f(a,b)}{M}[s] & = \Assign{f}{M}(1, 2) = 1+2 = 3.
\intertext{$f(f(a,b),a)$-এর মান গণনা করতে বিবেচনা করতে হবে}
\Value{f(f(a,b),a)}{M}[s] & = \Assign{f}{M}(\Value{f(a, b)}{M}[s],
\Value{a}{M}[s]) = \Assign{f}{M}(3, 1) = 3,
\intertext{কারণ $3+1 > 3$। $s(x) = 1$ এবং $\Value{x}{M}[s] =
  s(x)$ বলে আরও পাই}
\Value{f(f(a,b),x)}{M}[s] & = \Assign{f}{M}(\Value{f(a, b)}{M}[s],
\Value{x}{M}[s]) = \Assign{f}{M}(3, 1) = 3,
\end{align*}

পরমাণু !!{formula}~$R(t_1, t_2)$ পরিতৃপ্ত হয়, যদি তার যুক্তিগুলির
মানের টুপল, অর্থাৎ $\tuple{\Value{t_1}{M}[s],
  \Value{t_2}{M}[s]}$,~$\Assign{R}{M}$-এর !!a{element} হয়। তাই, যেমন,
$\Sat{M}{R(b,f(a,b))}[s]$, কারণ $\tuple{\Value{b}{M},
  \Value{f(a,b)}{M}} = \tuple{2, 3} \in \Assign{R}{M}$; কিন্তু
$\Sat/{M}{R(x, f(a,b))}[s]$, কারণ $\tuple{1, 3} \notin \Assign{R}{M}$।
% BN-SRC-093 TARGET-CORRECTION-BEGIN
\par\smallskip\noindent\textnormal{\small\textbf{উৎস-সংশোধন (BN-SRC-093):} হিমায়িত ইংরেজি উৎসে শেষ অ-সদস্যতার বিবৃতিতে সম্বন্ধের ব্যাখ্যা $\Assign{R}{M}$-এর পরে চলরাশি-আরোপ $[s]$ যুক্ত হয়েছে। আরোপটি পরমাণু সূত্রের পরিতৃপ্তি ও পদের মানে প্রযোজ্য, সম্বন্ধের স্থির ব্যাখ্যায় নয়; বাংলা পাঠে তাই $\Assign{R}{M}$ রাখা হয়েছে।} % BN-SRC-093 TARGET-CORRECTION-END

অপরমাণু সূত্র~$!A$ পরিতৃপ্ত কি না নির্ধারণ করতে প্রধান সংযোজকের জন্য
প্রযোজ্য আরোহী সংজ্ঞার অনুচ্ছেদটি ব্যবহার করতে হয়। যেমন,
$R(a, a) \lif (R(b,x) \lor R(x,b))$-এর প্রধান সংযোজক~$\lif$, এবং
\begin{align*}
  & \Sat{M}{R(a, a) \lif (R(b, x) \lor R(x, b))}[s] \text{ তখনই, যখন }\\
  & \qquad
  \Sat/{M}{R(a,a)}[s] \text{ অথবা } \Sat{M}{R(b, x) \lor R(x, b)}[s]
  \intertext{যেহেতু $\Sat{M}{R(a,a)}[s]$ (কারণ $\tuple{1,1} \in
    \Assign{R}{M}$), এখনও উত্তর নির্ধারণ করা যায় না; আগে বের করতে হবে
    $\Sat{M}{R(b, x) \lor R(x, b)}[s]$ কি না:}
  & \Sat{M}{R(b, x) \lor R(x, b)}[s] \text{ তখনই, যখন }\\
  & \qquad \Sat{M}{R(b,x)}[s] \text{ অথবা } \Sat{M}{R(x, b)}[s]
  \intertext{এটি সত্য, কারণ $\Sat{M}{R(x, b)}[s]$
    (যেহেতু $\tuple{1,2} \in \Assign{R}{M}$)।}
\end{align*}

মনে করুন, $s$-এর $x$-বিকল্প হলো এমন চলরাশি-আরোপ, যা $s$-এর থেকে
সর্বাধিক $x$-এর আরোপে আলাদা। $\Domain{M}$-এর প্রত্যেক !!{element}-এর
জন্য $s$-এর একটি $x$-বিকল্প আছে:
\begin{align*}
  s_1 & = \Subst{s}{1}{x}, &
  s_2 & = \Subst{s}{2}{x},\\
  s_3 & = \Subst{s}{3}{x}, &
  s_4 & = \Subst{s}{4}{x}.
\end{align*}
তাই যেমন, $s_2(x) = 2$ এবং $x$ ছাড়া সব চলরাশি~$y$-এর জন্য
$s_2(y) = s(y) = 1$। গঠন~$\Struct{M}$-এর জন্য এগুলিই $s$-এর সব
$x$-বিকল্প, কারণ $\Domain{M} = \{1, 2, 3, 4\}$। বিশেষভাবে লক্ষ করুন,
$s_1 = s$ ($s$ সব সময় নিজের একটি $x$-বিকল্প)।

\iftag{prvEx}{অস্তিত্বসূচকভাবে পরিমাণিত
  !!{formula}~$\lexists[x][!A(x)]$ পরিতৃপ্ত কি না নির্ধারণ করতে দেখতে
  হবে অন্তত একটি $m \in \Domain M$-এর জন্য
  $\Sat{M}{!A(x)}[\Subst{s}{m}{x}]$ কি না। তাই
  \[
  \Sat{M}{\lexists[x][(R(b,x) \lor R(x,b))]}[s],
  \]
  কারণ $\Sat{M}{R(b,x) \lor R(x, b)}[\Subst{s}{1}{x}]$
  ($\Subst{s}{3}{x}$-ও উদ্দেশ্য পূরণ করত)। কিন্তু
  \[
  \Sat/{M}{\lexists[x][(R(b,x) \land R(x,b))]}[s]
  \]
  কারণ যে $m \in \Domain{M}$-ই বাছি না কেন,
  $\Sat/{M}{R(b,x) \land R(x,b)}[\Subst{s}{m}{x}]$।}{}

\iftag{prvAll}{সার্বিকভাবে পরিমাণিত
  !!{formula}~$\lforall[x][!A(x)]$ পরিতৃপ্ত কি না নির্ধারণ করতে দেখতে
  হবে সব $m \in \Domain M$-এর জন্য
  $\Sat{M}{!A(x)}[\Subst{s}{m}{x}]$ কি না। তাই
  \[
  \Sat{M}{\lforall[x][(R(x,a) \lif R(a,x))]}[s],
  \]
  কারণ সব $m \in \Domain M$-এর জন্য
  $\Sat{M}{R(x,a) \lif R(a,x)}[\Subst{s}{m}{x}]$। $m = 1$ হলে
  $\Sat{M}{R(a,x)}[\Subst{s}{1}{x}]$, তাই অনুবর্তীটি সত্য; আর
  $m = 2$, $3$ ও~$4$ হলে $\Sat/{M}{R(x,a)}[\Subst{s}{m}{x}]$,
  তাই পূর্ববর্তীটি মিথ্যা। কিন্তু
  \[
  \Sat/{M}{\lforall[x][(R(a,x) \lif R(x,a))]}[s]
  \]
  কারণ $\Sat/{M}{R(a,x) \lif R(x,a)}[\Subst{s}{2}{x}]$ (যেহেতু
  $\Sat{M}{R(a, x)}[\Subst{s}{2}{x}]$ এবং $\Sat/{M}{R(x,
  a)}[\Subst{s}{2}{x}]$)।}{}

\iftag{defEx}{অস্তিত্বসূচকভাবে পরিমাণিত
  !!{formula}~$\lexists[x][!A(x)]$ পরিতৃপ্ত কি না নির্ধারণ করতে দেখতে
  হবে $\Sat{M}{\lnot \lforall[x][\lnot !A(x)]}[s]$ কি না। যেমন, আমাদের
  \[
  \Sat{M}{\lexists[x][(R(b,x) \lor R(x,b))]}[s].
  \]
  প্রথমে, $\Sat{M}{R(b,x) \lor R(x, b)}[\Subst{s}{1}{x}]$
  ($\Subst{s}{3}{x}$-ও উদ্দেশ্য পূরণ করত)। তাই
  $\Sat/{M}{\lnot(R(b,x) \lor R(x,b))}[\Subst{s}{1}{x}]$, ফলে
  $\Sat/{M}{\lforall[x][\lnot(R(b,x) \lor R(x,b))]}[s]$, এবং অতএব
  $\Sat{M}{\lnot\lforall[x][\lnot(R(b,x) \lor R(x,b))]}[s]$।
  অন্য দিকে,
  \[
  \Sat/{M}{\lexists[x][(R(b,x) \land R(x,b))]}[s].
  \]
  কারণ $\Sat{M}{\lforall[x][\lnot(R(b,x) \land R(x,b))]}[s]$;
  কোনো $m \in \Domain M$-এর জন্যই
  $\Sat{M}{R(b,x) \land R(x,b)}[\Subst{s}{m}{x}]$ নয়। এই উদাহরণগুলি
  থেকে অনুমান করা যায়, $\Sat{M}{\lexists[x][!A(x)]}[s]$ তখনই, যখন
  অন্তত একটি $m \in \Domain M$-এর জন্য
  $\Sat{M}{!A(x)}[\Subst{s}{m}{x}]$।}{}
% BN-SRC-094 TARGET-CORRECTION-BEGIN
\par\smallskip\noindent\textnormal{\small\textbf{উৎস-সংশোধন (BN-SRC-094):} হিমায়িত ইংরেজি উৎসের সংজ্ঞায়িত অস্তিত্বসূচকের উদাহরণে $\lnot$-এর পরে দুটি খোলা বন্ধনী কিন্তু একটি বন্ধ বন্ধনী ছিল—একবার সার্বিক সূত্রে এবং একবার তার নঞর্থকরণে। বাংলা পাঠে উভয় স্থানে সুগঠিত $\lnot(R(b,x) \lor R(x,b))$ রাখা হয়েছে।} % BN-SRC-094 TARGET-CORRECTION-END
% BN-SRC-095 TARGET-CORRECTION-BEGIN
\par\smallskip\noindent\textnormal{\small\textbf{উৎস-সংশোধন (BN-SRC-095):} হিমায়িত ইংরেজি উৎসে অপরিতৃপ্ত অস্তিত্বসূচক সূত্রটির $\Sat/$-যুক্তির ভিতরে সূত্রের পরে একটি অতিরিক্ত কমা রয়েছে। বাংলা পাঠে কমাটি বাদ দিয়ে $\lexists[x][(R(b,x) \land R(x,b))]$ রাখা হয়েছে।} % BN-SRC-095 TARGET-CORRECTION-END

\iftag{defAll}{সার্বিকভাবে পরিমাণিত
  !!{formula}~$\lforall[x][!A(x)]$ পরিতৃপ্ত কি না নির্ধারণ করতে দেখতে
  হবে $\Sat{M}{\lnot\lexists[x][\lnot !A(x)]}[s]$ কি না। যেমন,
  \[
  \Sat{M}{\lforall[x][(R(x,a) \lif R(a,x))]}[s],
  \]
  প্রথমে, সব $m \in \Domain M$-এর জন্য
  $\Sat{M}{R(x,a) \lif R(a,x)}[\Subst{s}{m}{x}]$
  ($\Sat{M}{R(a,x)}[\Subst{s}{1}{x}]$ এবং $m = 2$, $3$ বা~$4$ হলে
  $\Sat/{M}{R(x,a)}[\Subst{s}{m}{x}]$)। তাই এমন কোনো
  $m \in \Domain M$ নেই যার জন্য
  $\Sat{M}{\lnot(R(x,a) \lif R(a,x))}[\Subst{s}{m}{x}]$; ফলে
  $\Sat/{M}{\lexists[x][\lnot(R(x,a) \lif R(a,x))]}[s]$। অতএব
  $\Sat{M}{\lnot\lexists[x][\lnot(R(x,a) \lif R(a,x))]}[s]$। অন্য দিকে,
  \[
  \Sat/{M}{\lforall[x][(R(a,x) \lif R(x,a))]}[s],
  \]
  কারণ $\Sat/{M}{R(a,x) \lif R(x,a)}[\Subst{s}{2}{x}]$; তাই
  $\Sat{M}{\lexists[x][\lnot(R(a,x) \lif R(x,a))]}[s]$। এই
  উদাহরণগুলি থেকে অনুমান করা যায়, $\Sat{M}{\lforall[x][!A(x)]}[s]$
  তখনই, যখন প্রত্যেক $m \in \Domain M$-এর জন্য
  $\Sat{M}{!A(x)}[\Subst{s}{m}{x}]$।}{}
% BN-SRC-096 TARGET-CORRECTION-BEGIN
\par\smallskip\noindent\textnormal{\small\textbf{উৎস-সংশোধন (BN-SRC-096):} হিমায়িত ইংরেজি উৎসের সংজ্ঞায়িত সার্বিক পরিমাণসূচকের ব্যাখ্যায় $m=2,3,4$ ক্ষেত্রে অনুবর্তী $R(a,x)$-কে অপরিতৃপ্ত বলা হয়েছে; কিন্তু $R(a,2)$ সত্য এবং যুক্তিটি আসলে পূর্ববর্তী $R(x,a)$-এর অপরিতৃপ্তি ব্যবহার করে। বাংলা পাঠে তাই $R(x,a)$ রাখা হয়েছে।} % BN-SRC-096 TARGET-CORRECTION-END

আরও জটিল একটি ক্ষেত্র হিসেবে বিবেচনা করুন
\[
\lforall[x][(R(a,x) \lif \lexists[y][R(x,y)])].
\]
যেহেতু $\Sat/{M}{R(a,x)}[\Subst{s}{3}{x}]$ এবং
$\Sat/{M}{R(a,x)}[\Subst{s}{4}{x}]$, তাই শর্তবচনের অনুবর্তী নিয়ে
ভাবতে হয় এমন আকর্ষণীয় ক্ষেত্র কেবল $m = 1$ ও $m = 2$।
$\Sat{M}{\lexists[y][R(x,y)]}[\Subst{s}{1}{x}]$ কি সত্য? অন্তত একটি
$n \in \Domain M$-এর জন্য
$\Sat{M}{R(x,y)}[\Subst{\Subst{s}{1}{x}}{n}{y}]$ হলে তা সত্য। বস্তুত
$n = 1$ নিলে পাই $\Subst{\Subst{s}{1}{x}}{n}{y} =
\Subst{s}{1}{y} = s$। যেহেতু $s(x) = 1$, $s(y) = 1$ এবং
$\tuple{1,1} \in \Assign{R}{M}$, উত্তরটি হ্যাঁ।
% BN-SRC-097 TARGET-CORRECTION-BEGIN
\par\smallskip\noindent\textnormal{\small\textbf{উৎস-সংশোধন (BN-SRC-097):} হিমায়িত ইংরেজি উৎসে আকর্ষণীয় দ্বিতীয় ক্ষেত্রটি “$=2$” লেখা, ফলে চলরাশির নাম অনুপস্থিত। আগের ও পরের বাক্যে পরিমাণিত মান~$m$ নিয়ন্ত্রণ করে; বাংলা পাঠে $m=2$ রাখা হয়েছে।} % BN-SRC-097 TARGET-CORRECTION-END

$\Sat{M}{\lexists[y][R(x,y)]}[\Subst{s}{2}{x}]$ কি না নির্ধারণ করতে
!!{variable} আরোপগুলি $\Subst{\Subst{s}{2}{x}}{n}{y}$ দেখতে হবে।
এখানে $n = 1$ হলে এই আরোপটি~$s_2 = \Subst{s}{2}{x}$, যা $R(x,y)$-কে
পরিতৃপ্ত করে না ($s_2(x) = 2$, $s_2(y) = 1$ এবং
$\tuple{2,1}\notin \Assign{R}{M}$)। কিন্তু
$\Subst{\Subst{s}{2}{x}}{3}{y} = \Subst{s_2}{3}{y}$ বিবেচনা করুন।
$\tuple{2,3} \in \Assign{R}{M}$ বলে
$\Sat{M}{R(x,y)}[\Subst{s_2}{3}{y}]$; তাই
$\Sat{M}{\lexists[y][R(x,y)]}[s_2]$।

সুতরাং প্রত্যেক $m \in \Domain M$-এর জন্য হয়
$\Sat/{M}{R(a,x)}[\Subst{s}{m}{x}]$ ($m = 3$, $4$ হলে), নয়তো
$\Sat{M}{\lexists[y][R(x,y)]}[\Subst{s}{m}{x}]$ ($m = 1$, $2$ হলে);
অতএব
\[
\Sat{M}{\lforall[x][(R(a,x) \lif \lexists[y][R(x,y)])]}[s].
\]
অন্য দিকে,
\[
\Sat/{M}{\lexists[x][(R(a,x) \land \lforall[y][R(x,y)])]}[s].
\]
শুধু $m = 1$ ও $m = 2$-এর জন্যই
$\Sat{M}{R(a,x)}[\Subst{s}{m}{x}]$। কিন্তু এই দুই মানের ক্ষেত্রেই
এমন $n \in \Domain M$ আছে—$m=1$ হলে $n=4$ এবং $m=2$ হলে $n=1$—যাতে
$\Sat/{M}{R(x,y)}[\Subst{\Subst{s}{m}{x}}{n}{y}]$; তাই $m = 1$ ও
$m = 2$-এর জন্য $\Sat/{M}{\lforall[y][R(x,y)]}[\Subst{s}{m}{x}]$।
সব মিলিয়ে এমন কোনো $m \in \Domain M$ নেই যাতে
$\Sat{M}{R(a,x) \land \lforall[y][R(x,y)]}[\Subst{s}{m}{x}]$।% BN-SRC-098 TARGET-CORRECTION-BEGIN
\par\smallskip\noindent\textnormal{\small\textbf{উৎস-সংশোধন (BN-SRC-098):} হিমায়িত ইংরেজি উৎসে পরিমাণের শেষ সারাংশ “সব $n$”-এর জন্য বলা হলেও ব্যবহৃত চলরাশি ও ক্ষেত্রগুলি $m=1,2,3,4$; বাংলা পাঠে “প্রত্যেক $m$” রাখা হয়েছে।} % BN-SRC-098 TARGET-CORRECTION-END % BN-SRC-099 TARGET-CORRECTION-BEGIN
\par\smallskip\noindent\textnormal{\small\textbf{উৎস-সংশোধন (BN-SRC-099):} হিমায়িত ইংরেজি উৎসে $m=1$ ও $m=2$ উভয়ের জন্য $n=4$-কে $R(m,n)$-এর বিপরীত-উদাহরণ বলা হয়েছে। কিন্তু ঘোষিত সম্বন্ধে $\tuple{2,4}\in\Assign{R}{M}$। বাংলা পাঠে $m=1$-এর জন্য $n=4$ এবং $m=2$-এর জন্য $n=1$ নেওয়া হয়েছে; উভয় সংশ্লিষ্ট জোড়াই ঘোষিত সম্বন্ধের বাইরে।} % BN-SRC-099 TARGET-CORRECTION-END
\end{ex}

\iftag{defEx}{%
\begin{prop}\ollabel{prop:sat-ex}
$\Sat{M}{\lexists[x][!B(x)]}[s]$ তখনই, যখন $s$-এর এমন একটি
$x$-বিকল্প $s'$ আছে যাতে $\Sat{M}{!B(x)}[s']$।
\end{prop}

\begin{proof}
  অনুশীলনী।
\end{proof}
}{}

\tagprob{defEx}
\begin{prob}
  \olref[fol][syn][sat]{prop:sat-ex} প্রমাণ করুন।
\end{prob}
\tagendprob

\iftag{defAll}{%
\begin{prop}\ollabel{prop:sat-all}
$\Sat{M}{\lforall[x][!B(x)]}[s]$ তখনই, যখন $s$-এর প্রত্যেক
$x$-বিকল্প~$s'$-এর জন্য $\Sat{M}{!B(x)}[s']$।
\end{prop}

\begin{proof}
  অনুশীলনী।
\end{proof}
}{}

\tagprob{defAll}
\begin{prob}
  \olref[fol][syn][sat]{prop:sat-all} প্রমাণ করুন।
\end{prob}
\tagendprob

\begin{prob}
ধরি $\Lang L = \{c, f, A\}$, যেখানে একটি !!{constant}, একটি
এক-স্থানীয় !!{function} এবং একটি দ্বি-স্থানীয় !!{predicate} আছে;
আর !!{structure}~$\Struct{M}$ নিচের নিয়মে দেওয়া:
\begin{enumerate}
\item $\Domain M = \{1, 2, 3\}$
\item $\Assign{c}{M} = 3$
\item $\Assign{f}{M}(1) = 2, \Assign{f}{M}(2) = 3, \Assign{f}{M}(3) = 2$
\item $\Assign{A}{M} = \{\tuple{1, 2}, \tuple{2, 3}, \tuple{3, 3}\}$
\end{enumerate}
(ক) সব !!{variable}s~$v$-এর জন্য $s(v) = 1$ ধরি। নির্ধারণ করুন,
\[
\Sat{M}{\lexists[x][(A(f(z), c) \lif \lforall[y][(A(y, x) \lor A(f(y),
      x))])]}[s]
\]
কি না। কেন বা কেন নয়, ব্যাখ্যা করুন।

(খ) এমন একটি ভিন্ন গঠন ও !!{variable} আরোপ দিন যাতে !!{formula}-টি
পরিতৃপ্ত নয়।
\end{prob}

\end{document}
